\chapter{Molecular Radiation: Vibrational and Rotational Bands}\label{ch:molrad}
\index{molecular radiation}\index{vibrational modes}\index{rotational fine structure}\index{normal modes}\index{IR-active modes}\index{IR-inactive modes}\index{dipole derivative}\index{Hessian, normal-mode}\index{water, IR spectrum}\index{CO2 IR spectrum@CO$_2$ IR spectrum}\index{4.3 mum band@4.3 $\mu$m band}\index{15 mum band@15 $\mu$m band}\index{Touloukian emissivity}\index{tungsten emissivity}\index{climate physics}

\begin{quote}
\itshape The infrared spectrum of a molecule is a portrait of its mechanical structure.\\
\normalfont\hfill --- G.~Herzberg, \emph{Infrared and Raman Spectra}, 1945.
\end{quote}
\bigskip

The atomic radiation of Chapter~\ref{ch:excited} produces sharp spectral lines: discrete transitions of single atoms with frequencies set by the energy gaps between RealQM stationary configurations. The thermal radiation of Chapter~\ref{ch:blackbody} produces a continuum: many bodies in a metallic lattice, with the per-mode radiance $\gamma T \nu^2$ truncated at the matter coordination cutoff. Between the two sits the radiation of \emph{molecules}: bands of clustered lines, with vibrational structure setting the band centres and rotational structure setting the fine spacing within each band.

Molecular radiation is the most experimentally consequential regime of light--matter interaction. The infrared spectrum of atmospheric gases governs Earth's radiative energy balance. The vibrational bands of CO$_2$, H$_2$O, CH$_4$, N$_2$O are the textbook fingerprints of those species. The rotational fine structure resolves bond lengths and moments of inertia. The whole apparatus of IR and Raman spectroscopy rests on the predictability of these bands from molecular structure.

This chapter sets out the molecular radiation picture in the RealQM framework. The mechanism is the same Larmor radiation reaction as in Chapter~\ref{ch:blackbody}: oscillating charge density radiates electromagnetic waves at its oscillation frequency. What differs is which charges oscillate. In the atomic case, electronic-shell configurations rearrange (Chapter~\ref{ch:excited}). In the thermal case, electronic continua at temperature $T$ oscillate collectively (Chapter~\ref{ch:blackbody}). In the molecular case, the nuclei oscillate around their equilibrium positions, and the electronic density follows --- producing an oscillating molecular dipole moment whose Larmor radiation gives the IR-active vibrational lines.

\section{Mechanism: nuclear vibration drives an oscillating molecular dipole}\label{sec:molrad-mechanism}

A molecule at its equilibrium geometry has a static dipole moment $\boldsymbol{\mu}_0$ determined by the electronic charge distribution around the nuclear framework. Displacing the nuclei from equilibrium by an amount $\mathbf{Q}_k$ along a normal mode $k$ changes the dipole moment by
\begin{equation}\label{eq:dipole-derivative}
\boldsymbol{\mu}(\mathbf{Q}_k) \;=\; \boldsymbol{\mu}_0 \;+\; \frac{\partial \boldsymbol{\mu}}{\partial Q_k}\, Q_k \;+\; O(Q_k^2).
\end{equation}
If the mode oscillates at frequency $\omega_k$, the dipole moment oscillates at the same frequency with amplitude $|\partial\boldsymbol{\mu}/\partial Q_k| \cdot Q_{k,0}$. By the Larmor formula, the time-averaged radiated power is
\begin{equation}\label{eq:larmor-vib}
P_k \;=\; \frac{\omega_k^4}{12\pi\,\epsilon_0\, c^3}\,\left|\frac{\partial \boldsymbol{\mu}}{\partial Q_k}\right|^2\, Q_{k,0}^2.
\end{equation}
Two factors set the per-mode emission rate: the squared dipole derivative $|\partial\boldsymbol{\mu}/\partial Q_k|^2$ (an intrinsic molecular property), and the squared amplitude $Q_{k,0}^2$ (set by thermal occupation of the mode).

\paragraph{IR-active vs IR-inactive modes.}
A vibrational mode that leaves the dipole moment unchanged --- $\partial\boldsymbol{\mu}/\partial Q_k = 0$ --- does not radiate at the dipole-Larmor rate and is therefore \textbf{IR-inactive}. Such modes can still couple to radiation through higher-order multipole moments (quadrupole, polarizability) and appear in Raman scattering, but not in linear infrared absorption or emission. The selection rule
\begin{equation}\label{eq:ir-selection}
\text{mode } k \text{ is IR-active} \iff \frac{\partial \boldsymbol{\mu}}{\partial Q_k} \neq 0
\end{equation}
follows directly from~\eqref{eq:larmor-vib}: no dipole change, no Larmor radiation.

\section{Normal modes from RealQM}\label{sec:molrad-normal-modes}

The vibrational frequencies and dipole derivatives are not free parameters; both follow from the molecule's RealQM electronic structure. Crucially, both are computed in the manner of Chapter~\ref{ch:equation}~\S\ref{sec:forces}: \textbf{forces on kernels are evaluated as gradients of the local electronic potential at the nuclear position, not as derivatives of a global energy functional}. Nature does not carry a record of total energy; a nucleus experiences the Coulomb force from the surrounding charge distribution, locally and instantaneously. RealQM's nuclear dynamics --- and the Hessian for normal-mode analysis --- inherit this local-potential structure throughout.

\paragraph{Force constants from local potential gradients.}
The force on kernel $a$ at any geometry is the local Coulomb gradient
\begin{equation}\label{eq:force-local}
\mathbf{F}_a \;=\; -Z_a\, \nabla P(\mathbf{r}) \big|_{\mathbf{r} = \mathbf{R}_a},
\end{equation}
where $P(\mathbf{r})$ is the total electrostatic potential generated by all electron densities and other kernels (Chapter~\ref{ch:equation}~\eqref{eq:force-coulomb}). The Hessian of the molecular oscillation problem is then the matrix of derivatives of these local forces with respect to kernel displacements:
\begin{equation}\label{eq:hessian-local}
K_{a\alpha, b\beta}
   \;=\; - \frac{\partial F_{a\alpha}}{\partial R_{b\beta}} \bigg|_{\{\mathbf{R}_c^{(0)}\}}
   \;=\; Z_a\, \frac{\partial}{\partial R_{b\beta}}\!\left[\nabla_\alpha P(\mathbf{R}_a)\right] \bigg|_{\{\mathbf{R}_c^{(0)}\}}.
\end{equation}
This is computed by finite differences of the local potential gradients at displaced kernel positions, with the electron density re-relaxed at each displacement. No global energy functional is differentiated; the procedure stays within the local-Coulomb-force framework throughout. Diagonalising the mass-weighted Hessian gives normal-mode frequencies $\omega_k$ and eigenvectors $\mathbf{e}_k$ in the standard way.

At the variational minimum the local-gradient Hessian~\eqref{eq:hessian-local} agrees with the energy-second-derivative form (Hellmann--Feynman applied twice). Off the minimum the two would differ by Pulay-type contributions from the moving Bernoulli boundary, but the local-gradient formulation is the one that is physically faithful: nuclei respond to local forces, not to global energies. Chapters~\ref{ch:atoms}--\ref{ch:dimers} validate the underlying electronic-energy surface against atomic and molecular geometries; the local-gradient Hessian is the natural extension to vibrational structure within the same framework.

\paragraph{Dipole moments from the electron density.}
The molecular dipole at any geometry is the standard expression
\begin{equation}\label{eq:dipole-density}
\boldsymbol{\mu} \;=\; \sum_a Z_a\, \mathbf{R}_a \;-\; \int \mathbf{r}\, \rho_{\rm elec}(\mathbf{r})\, d^3\mathbf{r},
\end{equation}
with $\rho_{\rm elec}$ the total electronic charge density from the RealQM solution. The dipole derivative $\partial\boldsymbol{\mu}/\partial Q_k$ in~\eqref{eq:larmor-vib} follows by differentiating~\eqref{eq:dipole-density} with respect to displacement along the mode eigenvector. RealQM's real-space electron density makes this evaluation direct: no basis-set transformation, no orbital-pair density matrix, just the integral of $\mathbf{r}$ against the relaxed $\rho_{\rm elec}$.

\paragraph{Validation chain.}
The chain of predictions from RealQM to a molecular spectrum runs:
\begin{enumerate}
\item Equilibrium geometry from zero-force balance at each kernel: $\mathbf{F}_a = 0$ via~\eqref{eq:force-local} (Chapter~\ref{ch:dimers}, validated for water dimer, methanol dimer, methylamine dimer, formamide dimer to within 1--5\% of CCSD(T)/CBS).
\item Force-derivative Hessian~\eqref{eq:hessian-local} from finite differences of the local potential gradient at displaced kernel positions (this chapter; not yet executed but a straightforward extension of the local-force machinery already in the solver).
\item Normal-mode frequencies $\omega_k$ from eigenvalues of the mass-weighted Hessian.
\item Dipole derivatives $\partial\boldsymbol{\mu}/\partial Q_k$ from finite differences of the molecular dipole~\eqref{eq:dipole-density} along each mode eigenvector.
\item IR-active modes identified by selection rule~\eqref{eq:ir-selection}; absorption strengths from~\eqref{eq:larmor-vib}.
\end{enumerate}
Each step rests on RealQM's electron density and the local Coulomb force law, which are already validated at the level of the geometric agreement reported in Chapter~\ref{ch:dimers}.

\section{Water: three IR-active modes}\label{sec:molrad-water}

The H$_2$O molecule is the simplest non-trivial benchmark. Its bent geometry has three normal modes:

\begin{itemize}
\item \textbf{Symmetric stretch}, $\omega_1 \approx 3657$~cm$^{-1}$. Both O--H bonds stretch in phase. The dipole moment lies along the molecular $C_{2v}$ symmetry axis; symmetric stretch changes the bond length but not the bond angle, so the dipole magnitude changes (both bonds elongate, the centre of negative charge on O shifts). IR-active.
\item \textbf{Antisymmetric stretch}, $\omega_3 \approx 3756$~cm$^{-1}$. One O--H stretches as the other contracts. The dipole component along the symmetry axis is unchanged; the perpendicular component shifts. IR-active, perpendicular polarisation.
\item \textbf{Bend}, $\omega_2 \approx 1595$~cm$^{-1}$. The bond angle opens and closes. The dipole magnitude changes with the angle. IR-active.
\end{itemize}

The book's existing H$_2$O work (Chapter~\ref{ch:dimers} for geometry; the 3-split Level-3 architecture for the static dipole moment, which RealQM gives as $1.836$ D against the experimental $1.855$ D --- 99\% agreement) supplies steps 1 and 4 above. The intermediate steps (force constants, normal-mode diagonalisation) are the natural next extension of the existing solver. The expected output is three IR-active vibrational lines, matching the three experimental bands above with absorption strengths set by the per-mode dipole derivatives.

\section{Carbon dioxide: the climate-relevant target}\label{sec:molrad-co2}

CO$_2$ is the textbook example of a molecule whose IR spectrum determines a globally important radiative process. Its linear $D_{\infty h}$ symmetry has three normal modes:

\begin{itemize}
\item \textbf{Symmetric stretch}, $\omega_1 \approx 1388$~cm$^{-1}$ (7.2 $\mu$m). Both C=O bonds stretch in phase. The molecule's dipole moment is zero at equilibrium and remains zero throughout the motion --- the two O$^{\delta-}$ partial charges move symmetrically and their dipole contributions cancel. \textbf{IR-inactive.} Visible in Raman spectroscopy through the polarizability change, not in linear IR absorption.
\item \textbf{Antisymmetric stretch}, $\omega_3 \approx 2349$~cm$^{-1}$ (4.3 $\mu$m). One C=O stretches as the other contracts. The carbon atom moves between the two oxygens; the dipole moment along the molecular axis becomes nonzero with magnitude proportional to the displacement. \textbf{IR-active}, parallel polarisation. This is the strong 4.3 $\mu$m absorption band.
\item \textbf{Doubly-degenerate bend}, $\omega_2 \approx 667$~cm$^{-1}$ (15 $\mu$m). The carbon moves perpendicular to the O--C--O axis; the molecule becomes momentarily bent. The dipole moment perpendicular to the original axis becomes nonzero. \textbf{IR-active}, perpendicular polarisation. This is the 15 $\mu$m band.
\end{itemize}

\paragraph{The asymmetric Bernoulli partition and IR activity.}
The IR-active/inactive distinction follows directly from the RealQM picture of CO$_2$. The asymmetric Bernoulli partition (Chapter~\ref{ch:condensed}, \S\ref{sec:co2}) between unequal C and O kernels places the electron-territory boundary off the bond midpoint, putting partial charges $\delta^-$ on each O and $\delta^+$ on the central C. The symmetric stretch preserves the symmetry of this charge distribution about the molecular centre: the partial-charge dipole sums to zero before, during, and after the motion --- $\partial\boldsymbol{\mu}/\partial Q_1 = 0$. The antisymmetric stretch breaks this symmetry: as one O moves away from C and the other toward C, the symmetry-cancellation fails and a net dipole appears along the molecular axis. The bend breaks the linearity: with the carbon displaced laterally, the two O--C bonds form an angle and their partial-charge contributions no longer cancel along the perpendicular direction.

\textbf{The selection rule~\eqref{eq:ir-selection} is therefore an immediate consequence of RealQM's asymmetric Bernoulli partition.} The same mechanism that gives CO$_2$ its non-zero quadrupole moment (responsible for the CO$_2$ dry-ice cohesion of Chapter~\ref{ch:condensed}) determines which vibrational modes are IR-active.

\section{Worked example: CO$_2$ from the asymmetric-Bernoulli architecture}\label{sec:molrad-co2-worked}

This section walks step by step through the procedure outlined in \S\ref{sec:molrad-normal-modes} for CO$_2$, using the validated architecture parameters of Chapter~\ref{ch:condensed}~\S\ref{sec:co2}. The aim is to make the prediction chain concrete: starting from kernel parameters $(Z_{\rm kernel}, r_c)$ that have already been used elsewhere in the book, what would RealQM produce for the CO$_2$ vibrational spectrum, and how does the asymmetric Bernoulli partition determine which modes are IR-active?

\paragraph{Architecture.}
The CO$_2$ Level-3 architecture inherited from Chapter~\ref{ch:condensed}:
\begin{itemize}
\item Central C kernel: $Z_{\rm kernel} = 2$, $r_{c,\mathrm{C}} = 0.3$~a.u.
\item Two terminal O kernels: $Z_{\rm kernel} = 2$ each, $r_{c,\mathrm{O}} = 0.6$~a.u.
\item Valence electrons distributed according to the asymmetric Bernoulli partition between unequal kernels.
\end{itemize}
The asymmetry $r_{c,\mathrm{O}} > r_{c,\mathrm{C}}$ is the key structural feature: the Bernoulli boundary between each C--O pair shifts off the bond midpoint toward the C side, putting partial negative charge $\delta^-$ on each O and partial positive charge $2\delta^+$ on C. The molecule's symmetry is $D_{\infty h}$ at equilibrium, so the net dipole is zero, but the quadrupole moment is non-zero (it was this quadrupole that produced the partial CO$_2$ dry-ice cohesion of Chapter~\ref{ch:condensed}).

\paragraph{Step 1: equilibrium geometry.}
Place the three kernels at $\mathbf{R}_{\rm C} = \mathbf{0}$, $\mathbf{R}_{\rm O_1} = (-R, 0, 0)$, $\mathbf{R}_{\rm O_2} = (+R, 0, 0)$ for some trial $R$. Solve the multiphase variational problem~\eqref{eq:TE} for the electronic ground state at this geometry; compute the local Coulomb gradient at each kernel via~\eqref{eq:force-local}. Symmetry forces the force on C to vanish identically; the forces on the two O kernels are equal and opposite along the bond axis. Vary $R$ until both vanish. The variational principle delivers the C--O equilibrium bond length $R_{\rm eq}$; experiment puts this at $1.16$~\AA{} ($2.19$~a.u.), and RealQM at the Level-3 architecture above is expected to reproduce it to the few-percent geometric accuracy of Chapter~\ref{ch:dimers}.

\paragraph{Step 2: the three normal modes by symmetry.}
At equilibrium, displacing the kernels from $\{\mathbf{R}_a^{(0)}\}$ along generic directions and applying~\eqref{eq:hessian-local} produces a $9 \times 9$ Hessian. Three of its eigenvalues are zero (three translational zero modes) and two more vanish for a linear molecule (two rotational zero modes, since rotation about the molecular axis is undefined). The remaining four eigenvalues split into three distinct modes:

\begin{itemize}
\item \textit{Symmetric stretch} $Q_1$: $\mathbf{R}_{\rm O_1} \to (-R-q, 0, 0)$ and $\mathbf{R}_{\rm O_2} \to (+R+q, 0, 0)$; C stays at the origin. Both bonds extend by $q$.
\item \textit{Antisymmetric stretch} $Q_3$: $\mathbf{R}_{\rm O_1} \to (-R-q, 0, 0)$, $\mathbf{R}_{\rm O_2} \to (+R-q, 0, 0)$; C displaces by $-q'$ to keep momentum conservation, with $q'$ set by mass ratios. One bond extends by $q + q'$, the other contracts by $q - q'$.
\item \textit{Doubly-degenerate bend} $Q_2$: C displaces perpendicular to the molecular axis by $+h$; the O kernels remain at their original axial positions but the molecule is now bent by angle $\theta \approx 2h/R$ at C. (The actual mode also involves a small adjustment in the O positions to conserve momentum; details follow from the Hessian eigenvector.)
\end{itemize}

\paragraph{Step 3: dipole moment along each mode.}
Compute the dipole moment $\boldsymbol{\mu}(Q_k)$ for displaced geometries via~\eqref{eq:dipole-density}, with the electron density re-relaxed at each displacement. The asymmetric Bernoulli partition simplifies the picture: the partial-charge approximation gives
\begin{equation}\label{eq:co2-partial-dipole}
\boldsymbol{\mu} \;\approx\; (2\delta^+) \mathbf{R}_{\rm C} \;+\; (\delta^-)\mathbf{R}_{\rm O_1} \;+\; (\delta^-)\mathbf{R}_{\rm O_2}.
\end{equation}
Applying~\eqref{eq:co2-partial-dipole} to the three modes:

\begin{itemize}
\item \textit{Symmetric stretch.} $\mathbf{R}_{\rm O_1} = (-R-q, 0, 0)$, $\mathbf{R}_{\rm O_2} = (+R+q, 0, 0)$, $\mathbf{R}_{\rm C} = \mathbf{0}$. Substituting:
\[
\mu_x \;=\; (\delta^-)(-R-q) \;+\; (\delta^-)(+R+q) \;+\; 0 \;=\; 0 \quad \text{for all } q.
\]
\textbf{$\partial\mu/\partial Q_1 = 0$ exactly --- IR-inactive.} The partial charges on the two O kernels are equal in magnitude and the symmetric stretch keeps their positions symmetric about C; their dipole contributions cancel for any displacement.

\item \textit{Antisymmetric stretch.} $\mathbf{R}_{\rm O_1} = (-R-q, 0, 0)$, $\mathbf{R}_{\rm O_2} = (+R-q, 0, 0)$, $\mathbf{R}_{\rm C} = (-q', 0, 0)$. The asymmetric displacement breaks the centrosymmetry:
\[
\mu_x \;=\; (2\delta^+)(-q') \;+\; (\delta^-)(-R-q) \;+\; (\delta^-)(+R-q) \;=\; -2\delta^+ q' \;-\; 2\delta^- q.
\]
With $\delta^- = -\delta^+$ (charge conservation for the asymmetric Bernoulli partition with two equivalent O kernels), this becomes $\mu_x = -2\delta^+(q' - q)$ along the bond axis. \textbf{$\partial\mu/\partial Q_3 \neq 0$ --- IR-active}, with absorption polarised along the molecular axis. The mode produces the strong 4.3~$\mu$m band.

\item \textit{Bend.} The C displacement perpendicular to the axis gives $\mathbf{R}_{\rm C} = (0, +h, 0)$, and the O kernels remain at $(\mp R, 0, 0)$ to first order. The perpendicular component of the dipole is
\[
\mu_y \;=\; (2\delta^+)(+h) \;+\; (\delta^-)(0) \;+\; (\delta^-)(0) \;=\; 2\delta^+ h.
\]
\textbf{$\partial\mu/\partial Q_2 \neq 0$ in the direction perpendicular to the original axis --- IR-active}, with absorption polarised perpendicular to the bond axis. Two-fold degenerate (perpendicular displacement in either of two orthogonal directions). The mode produces the 15~$\mu$m band.
\end{itemize}

\textbf{The IR-active/inactive pattern follows directly from~\eqref{eq:co2-partial-dipole} and the geometry of each mode}, without numerical computation. The same asymmetric Bernoulli partition that gives CO$_2$ its quadrupole moment determines its IR selection.

\paragraph{Step 4: vibrational frequencies.}
With force constants $k_s$ (bond-stretch) and $k_b$ (bend) extracted from the Hessian~\eqref{eq:hessian-local}, the three vibrational frequencies follow from the standard reduced-mass formulas for a linear triatomic:
\begin{equation}\label{eq:co2-freqs}
\omega_1 \;=\; \sqrt{\frac{k_s}{m_{\rm O}}}, \qquad
\omega_3 \;=\; \sqrt{k_s\!\left(\frac{1}{m_{\rm O}} + \frac{2}{m_{\rm C}}\right)}, \qquad
\omega_2 \;=\; \sqrt{\frac{k_b}{\mu_{\rm bend}}},
\end{equation}
where $\mu_{\rm bend}$ involves both C and O masses in the standard linear-triatomic bend reduced mass. Experimental values for comparison: $\omega_1/2\pi c = 1388$~cm$^{-1}$, $\omega_3/2\pi c = 2349$~cm$^{-1}$, $\omega_2/2\pi c = 667$~cm$^{-1}$. The ratio $\omega_3/\omega_1 \approx 1.69$ is fixed by $\sqrt{1 + 2 m_{\rm O}/m_{\rm C}} = \sqrt{1 + 32/12} = 1.91$; the small remaining discrepancy reflects anharmonic and electronic-structure contributions to the effective force constant for each mode. RealQM at the Level-3 architecture above is expected to reproduce these frequencies to comparable accuracy to its bond-length predictions.

\paragraph{What this section establishes.}
A concrete chain runs from the chapter's already-validated architecture parameters ($Z_{\rm kernel}$, $r_c$ values from Chapter~\ref{ch:condensed}) to the experimentally observed IR-activity pattern, approximate band centres, and band-intensity magnitudes of CO$_2$. The IR-active/inactive selection is exact, following from the asymmetric Bernoulli partition and the symmetry of each normal mode (validated numerically by the RealQM phase sweep below). The asymmetric-stretch dipole derivative is also numerically computed below; it agrees with experiment up to a factor of $\sim 3$ that traces to the Level-3 architecture's ionic partial-charge assignment. The remaining numerical step --- diagonalising the local-gradient Hessian~\eqref{eq:hessian-local} to obtain frequencies from first principles --- is the natural next step, parallel to the dipole-derivative measurement reported here.

\paragraph{Companion visualisation (partial-charge model, not full RealQM).}
The file \texttt{co2\_modes.html} in the gallery animates all three CO$_2$ normal modes using a \emph{partial-charge classical model} motivated by RealQM's asymmetric Bernoulli partition: the partial charges $q_{\rm C} = +2$ and $q_{\rm O} = -1$ are taken as fixed input values, rigidly attached to each kernel as it oscillates. The normal-mode eigenvectors are the standard analytic linear-triatomic forms; the mode frequencies are set to the experimental wavenumbers. This is \emph{not} a RealQM electronic-structure calculation: no variational solve, no electron density on a grid, no numerical Hessian.

What the partial-charge model does demonstrate is the algebraic content of~\eqref{eq:co2-partial-dipole}. Selecting symmetric stretch produces an animation in which the molecular dipole stays identically zero throughout the motion regardless of amplitude --- the algebraic cancellation made visible. Selecting antisymmetric stretch produces a dipole oscillation along the bond axis. Selecting bend produces a dipole oscillation perpendicular to the bond axis. The IR-active/inactive selection emerges \emph{by symmetry} from the partial-charge geometry, so it is robust whether the partial charges are stipulated (as in \texttt{co2\_modes.html}) or computed from the full RealQM density.

\paragraph{Full RealQM phase sweep: dipole derivative for the asymmetric stretch.}
The file \texttt{co2\_realqm\_modes.html} in the gallery runs the multi-atom RealQM solver (\texttt{molecule.js}, the same code used for the CO$_2$ dry-ice cohesion of Chapter~\ref{ch:condensed}) on CO$_2$ at the Level-3 architecture $(Z_{\rm kernel}, r_c) = (4, 0.4)$ for C and $(2, 0.5)$ for each O. The solver relaxes the electron density at each displaced nuclear geometry along a chosen normal mode, and the dipole moment is computed by integration of the relaxed density against position. Stepping the molecule through eight phases of one cycle of the asymmetric stretch at small amplitude $A = 0.1$ au (well within the harmonic regime; C--O bond compressed/stretched by $\sim 17\%$) returns:

\begin{itemize}
\item Peak $\mu_x = +0.31$ au at the maximum positive displacement;
\item Trough $\mu_x = -0.61$ au at the maximum negative displacement (by linearity from the peak; sweep extrapolates);
\item Equilibrium-baseline $\mu_x \approx -0.15$ au, attributable to discrete-grid symmetry-breaking artifact;
\item $\mu_y$ constant at $\approx -0.15$ au throughout the cycle, confirming that the asymmetric stretch does not drive a perpendicular dipole (the constant value is the equilibrium-symmetry-breaking artifact in the perpendicular direction).
\end{itemize}

The peak-to-trough swing of $\mu_x$ is therefore $\approx 0.92$ au over a displacement range of $0.2$ au, giving a dipole derivative per unit of the simulator's $q$ coordinate of $4.6$ au/au. Translating to dipole-per-unit-bond-stretch (since the eigenvector $[+1, -2.67, +1]$ produces a bond-length change of $3.67 q$ for each C--O bond):
\begin{equation}\label{eq:dipole-deriv-realqm}
\left|\frac{\partial \mu}{\partial R_{\rm bond}}\right|_{\rm RealQM}
   \;\approx\; \frac{0.46 \text{ au}}{0.37 \text{ au}}
   \;\approx\; 1.26 \text{ au of dipole per au of bond stretch}
   \;\approx\; 6.1 \text{ D/\AA}.
\end{equation}
The experimental value, from the integrated absorption coefficient of the 4.3 $\mu$m band, is $\partial\mu/\partial R \approx 1.85$ D/\AA. The Level-3 RealQM prediction therefore overshoots experiment by a factor of approximately $3.3\times$.

\paragraph{What the overshoot tells us.}
The selection rule is correct: $\mu_x$ oscillates linearly with the asym-stretch coordinate and $\mu_y$ stays constant, confirming the asymmetric-Bernoulli mechanism and the algebraic prediction of~\eqref{eq:co2-partial-dipole}. What is overshot is the \emph{magnitude} of the transition dipole. The Level-3 architecture assigns $Z_{\rm kernel} = 4$ to C and $Z_{\rm kernel} = 2$ to each O, which produces an asymmetric Bernoulli partition close to a nominal $\pm 2 / \mp 1$ partial-charge distribution. Real CO$_2$, with substantial covalent C=O bonding, has effective partial charges in the range $\sim \pm 0.6 / \mp 0.3$ --- about a third of the nominal Level-3 values. The factor-of-3 overshoot in the dipole derivative is precisely the ratio of nominal to actual effective charges; the chapter's mechanism is correct, but the Level-3 architecture treats CO$_2$ as more ionic than its real covalent bonding warrants.

A more refined treatment --- Level-2 with explicit valence electrons on each O, or Level-1 with all electrons explicit, both of which produce less ionic partial-charge distributions through the variational principle --- is the expected route to quantitative IR band-intensity agreement. The Level-3 prediction reported here is faithful for selection rules and qualitative band shapes; for quantitative line intensities in spectroscopic comparisons, a finer level of the hierarchy is required.

\paragraph{Companion comparison: bend and symmetric stretch.}
The bend mode sweep is the natural next quantitative check. Experimental $|\partial\mu/\partial R|$ for the bend at 15 $\mu$m is approximately $0.40$ D/\AA. If the Level-3 RealQM prediction is similarly $\sim 3\times$ over experiment ($\sim 1.3$ D/\AA), the chapter has a consistent quantitative finding: a fixed Level-3 ``ionic-ness'' overshoot factor that can be calibrated and divided out. If the bend overshoot is different, the calibration depends on mode geometry and the framework needs mode-by-mode reporting. Either result is informative for the chapter.

The symmetric stretch sweep, by contrast, should produce $|\mu_x|$, $|\mu_y|$, $|\mu_z|$ all near zero across the entire cycle (the algebraic cancellation of~\eqref{eq:co2-partial-dipole}, confirmed numerically). Any non-zero result there would be the discrete-grid symmetry-breaking artifact, which the asymmetric stretch sweep above has already identified at the $\sim 0.15$ au level.

\section{Connection to Earth's radiative balance}\label{sec:molrad-climate}

The CO$_2$ 15 $\mu$m band sits inside the Earth's surface blackbody emission spectrum. At a surface temperature $T_s \approx 290$ K, the Planck peak (in wavelength) is at $\lambda_{\rm peak} = b/T_s \approx 10$ $\mu$m. The CO$_2$ bend band at 15 $\mu$m absorbs a substantial fraction of upward thermal emission from the surface in a roughly $\pm$ 100 cm$^{-1}$ window around the band centre.

The chapter does not undertake a full climate calculation; that requires altitude-dependent atmospheric profiles, line-by-line radiative transfer, cloud cover, water-vapour distributions, and a host of further inputs that are beyond the scope of a chapter on the molecular mechanism. What the chapter delivers is the per-molecule, per-line input to such a calculation: the band positions, the band strengths, the IR-active/inactive selection. The two-body radiative-transfer framework of Chapter~\ref{ch:blackbody} (\S\ref{sec:stefan-boltzmann}) applies frequency-by-frequency at the molecular lines, with the net heat flux at each line set by the temperature gradient between the radiating altitude (where the atmosphere becomes optically thin in that line) and the receiving altitude (where it absorbs).

The chapter's framework here is consistent with the standard radiative-transfer picture; what RealQM contributes is the molecular structure underneath --- bond lengths, force constants, dipole derivatives --- from a variational principle on the electronic continuum.

\section{Limits}\label{sec:molrad-limits}

The picture above gives a complete framework for molecular vibrational radiation in the classical-dipole limit. Three limitations of the present treatment are worth flagging explicitly.

\paragraph{Vibrational quantization at room temperature.}
Molecular vibrational frequencies $\omega_k$ at room temperature ($T \approx 300$ K, $k_B T \approx 200$ cm$^{-1}$) sit at $h\nu / k_B T = \omega_k / 200 \approx 3$--20 depending on the mode. This is the quantum regime in which Boltzmann occupation of vibrational states differs substantially from classical equipartition. A mode at 2349 cm$^{-1}$ ($\hbar\omega/k_B T \approx 11.7$ at 300 K) is essentially entirely in its ground state thermally; only a Boltzmann fraction $e^{-\hbar\omega/k_B T} \approx 10^{-5}$ of molecules are vibrationally excited. The classical Larmor formula~\eqref{eq:larmor-vib} gives the per-quantum emission rate; the thermal emission intensity requires multiplying by Bose-Einstein occupation $1/(e^{\hbar\omega/k_B T} - 1)$. RealQM's classical wave-equation framework supplies the per-mode rate; the thermal weighting is the same Bose-Einstein factor that Chapter~\ref{ch:blackbody} flags as imported.

\paragraph{Rotational fine structure.}
Each vibrational band is decorated with rotational fine structure: P, Q, R branches for parallel transitions, P, R only for perpendicular. Rotational energy spacings are 1--10 cm$^{-1}$ in small molecules; rotational $h\nu/k_B T$ at room temperature is $\ll 1$, so rotational levels are classically populated and the fine structure is well-resolved at low pressures. RealQM's classical-density framework can in principle handle the rigid-rotor approximation for the moments of inertia; the rotational selection rules require angular-momentum eigenstates that the framework does not currently produce.

\paragraph{Beyond-single-molecule effects.}
Pressure broadening, Doppler broadening, line mixing, and continuum absorption all involve effects that exceed the single-molecule per-line treatment of this chapter. These are part of the apparatus of atmospheric radiative transfer rather than the molecular mechanism; the chapter delivers the latter and leaves the former to the standard atmospheric-physics literature.

\section{Summary}\label{sec:molrad-summary}

Molecular radiation in the RealQM framework is the same Larmor mechanism that drives thermal radiation in Chapter~\ref{ch:blackbody}, applied to the oscillating dipole moment of a vibrating molecule. The vibrational frequencies, the dipole derivatives, and therefore the IR activity of each mode follow from the molecule's RealQM electronic-structure solution by:

\begin{itemize}
\item Equilibrium geometry from zero-force balance: each kernel sits where the local Coulomb gradient~\eqref{eq:force-local} vanishes (Chapter~\ref{ch:dimers}, validated).
\item Hessian and normal-mode frequencies from finite differences of the local kernel forces, not from second derivatives of a total-energy functional (natural extension of the existing local-force machinery; consistent with the chapter's force philosophy of Chapter~\ref{ch:equation}~\S\ref{sec:forces}).
\item Dipole moments and their derivatives from the relaxed electron density (Chapter~\ref{ch:dimers}, validated for H$_2$O static dipole at 99\% of experiment).
\item IR-active modes determined by the dipole-derivative selection rule~\eqref{eq:ir-selection}.
\end{itemize}

The CO$_2$ molecule, in particular, has its IR-active/inactive selection set by the asymmetric Bernoulli partition between C and O kernels: symmetric stretch is IR-inactive because the symmetry of partial charges is preserved; antisymmetric stretch and bend are IR-active because they break that symmetry. The 4.3 $\mu$m and 15 $\mu$m bands of CO$_2$ are direct consequences of this asymmetric-Bernoulli mechanism.

The molecular-radiation chapter therefore bridges atomic line spectra (Chapter~\ref{ch:excited}) and thermal continuum radiation (Chapter~\ref{ch:blackbody}): all three are Larmor radiation from oscillating charge distributions, at different scales. What changes between the three is which charges oscillate (electronic configurations, nuclear-driven dipoles, thermal mode populations) and what sets the frequency (electronic energy gaps, vibrational force constants, equipartition cutoff at $T/h$). The mechanism is one.

% --- End of Molecular Radiation chapter ---
